%% pst-numeric.tex: decimal numbers with any number of digits, matrices,
%% complex numbers and numerical methods in PostScript, for PSTricks
%%
%% Copyright 2012, 2026 Buddy Ledger
%%
%% This work may be distributed and/or modified under the
%% conditions of the LaTeX Project Public License, either version 1.3c
%% of this license or (at your option) any later version.
%% The latest version of this license is in
%%   https://www.latex-project.org/lppl.txt
%% and version 1.3c or later is part of all distributions of LaTeX
%% version 2008 or later.
%%
%% This work has the LPPL maintenance status `author-maintained'.
%%
%% The Current Maintainer of this work is Buddy Ledger.
%%
%% This work consists of the files listed in README.md.
%%
\csname PSTnumericLoaded\endcsname
\let\PSTnumericLoaded\endinput
\ifx\PSTricksLoaded\endinput\else\input pstricks.tex\fi
\ifx\PSTplotLoaded\endinput\else\input pst-plot.tex\fi
%
\edef\PstAtCode{\the\catcode`\@}\catcode`\@=11\relax
%
\def\fileversion{1.4}
\def\filedate{2026/11/22}
\message{`PST-numeric' v\fileversion, \filedate\space (BL)}
%
% The parts to load, a comma-separated list without spaces. Define
% \pstNumericParts before loading, or (in LaTeX) use the package options,
% which set it. The default is every part.
%   bignum   decimal numbers with any number of digits (bignumlibrary.ps)
%   trig     CORDIC trigonometry on reals (Trig.ps)
%   matrix   matrices of reals or bignums (Matrix.ps)
%   complex  complex numbers (Complex.ps), needs matrix
%   calc     roots, integrals, polynomials, interpolation, statistics,
%            regression, differential equations, Fourier transform
%            (Numeric.ps), needs matrix, complex and special
%   special  gamma, error and Bessel functions, distributions
%            (Special.ps), needs bignum
%   tensor   tensors of any rank (Tensor.ps), needs matrix
\ifx\pstNumericParts\pstnum@undefined
  \def\pstNumericParts{bignum,trig,matrix,complex,calc,special,tensor}\fi
%
\def\pstnum@empty{}
\def\pstnum@relax{\relax}
\def\pstnum@gobble#1{}
\def\pstnum@known{,bignum,trig,matrix,complex,calc,special,tensor,}
\def\pstnum@want#1{\expandafter\let\csname pstnum@want@#1\endcsname\pstnum@empty}
\def\pstnum@check#1#2{% #1: what follows #2 in the known list, empty if #2 is unknown
  \def\pstnum@tmp{#1}%
  \ifx\pstnum@tmp\pstnum@empty
    \errmessage{pst-numeric: unknown part `#2' (use bignum, trig, matrix, complex, calc, special or tensor)}%
  \else\pstnum@want{#2}\fi}
\def\pstnum@mark#1,{%
  \def\pstnum@tmp{#1}%
  \ifx\pstnum@tmp\pstnum@relax\else
    \ifx\pstnum@tmp\pstnum@empty\else
      \def\pstnum@try##1,#1,##2\pstnum@end{\pstnum@check{##2}{#1}}%
      \expandafter\pstnum@try\pstnum@known#1,\pstnum@end
    \fi
    \expandafter\pstnum@mark
  \fi}
\expandafter\pstnum@mark\pstNumericParts,\relax,
% parts that other parts need
\ifcsname pstnum@want@calc\endcsname\pstnum@want{complex}\pstnum@want{matrix}\pstnum@want{special}\fi
\ifcsname pstnum@want@tensor\endcsname\pstnum@want{matrix}\fi
\ifcsname pstnum@want@complex\endcsname\pstnum@want{matrix}\fi
\ifcsname pstnum@want@special\endcsname\pstnum@want{bignum}\fi
%
\pstheader{pst-numeric.pro}% the helpers of this interface
\def\pstnum@load#1{%
  \ifcsname pstnum@want@#1\endcsname
    \expandafter\pstheader\else\expandafter\pstnum@gobble\fi
  {pst-numeric-#1.pro}}
\pstnum@load{bignum}
\pstnum@load{trig}
\pstnum@load{matrix}
\pstnum@load{complex}
\pstnum@load{calc}
\pstnum@load{special}
\pstnum@load{tensor}
%
% The dictionary that holds every procedure of the library.
\def\pstNumericDict{tx@NumericDict}
% PostScript code run with the library's procedures, for use inside any
% PostScript argument of PSTricks:  \psplot{0}{1}{\pstNumeric{x 2 mul}}
\def\pstNumeric#1{ tx@NumericDict begin #1 end }
%
% ---------------------------------------------------------------- presentation
% Keys of the family pst-numeric, for \psset and the optional arguments:
%   numFormat   raw (as computed: all digits of a bignum), auto (like %g),
%               fixed (numDigits decimals), sci (numDigits significant digits)
%   numDigits   significant digits (auto, sci) or decimals (fixed)
%   numTrim     drop trailing zeros
%   numDecimal  the decimal separator, a single character
%   sciStyle    e (1.5e-3) or times (1.5 x 10 with a raised -3)
%   numFont, numFontSize   PostScript font name and size in pt
%   matrixFont  the font of \psNumericMatrix and \psNumericRegression
%   matrixDelim bracket, paren, bar, dbar or none
%   matrixAlign decimal, right, center or left
%   colSep, rowSep   column gap and row pitch, in units of the font size
%   vector      column or row (how a list of numbers is drawn)
%   label       drawn before the result followed by " = " (^{} and _{} work)
%   tensorLabels  slice labels such as (1,:,:) above a tensor's slices
%   regShow     the lines of \psNumericRegression, from eq params coef r2
%               adjr2 s n dof tval sse
%   regPlot     what \psNumericRegPlot draws: fit, ci, pi, data, residuals
%               or stdresiduals
\pst@addfams{pst-numeric}
\define@key[psset]{pst-numeric}{numFormat}{\def\psk@numFormat{#1}}
\define@key[psset]{pst-numeric}{numDigits}{\pst@getint{#1}\psk@numDigits}
\define@boolkey[psset]{pst-numeric}[Pst@]{numTrim}[true]{}
\define@key[psset]{pst-numeric}{numDecimal}{\def\psk@numDecimal{#1}}
\define@key[psset]{pst-numeric}{sciStyle}{\def\psk@sciStyle{#1}}
\define@key[psset]{pst-numeric}{numFont}{\def\psNumericFont{#1}}
\define@key[psset]{pst-numeric}{numFontSize}{\def\psNumericFontSize{#1}}
\define@key[psset]{pst-numeric}{matrixFont}{\def\psNumericMatrixFont{#1}}
\define@key[psset]{pst-numeric}{matrixDelim}{\def\psk@matrixDelim{#1}}
\define@key[psset]{pst-numeric}{matrixAlign}{\def\psk@matrixAlign{#1}}
\define@key[psset]{pst-numeric}{colSep}{\def\psk@colSep{#1}}
\define@key[psset]{pst-numeric}{rowSep}{\def\psk@rowSep{#1}}
\define@key[psset]{pst-numeric}{vector}{\def\psk@vector{#1}}
\define@key[psset]{pst-numeric}{label}{\def\psk@numLabel{#1}}
\define@boolkey[psset]{pst-numeric}[Pst@]{tensorLabels}[true]{}
\define@key[psset]{pst-numeric}{regShow}{\def\psk@regShow{#1}}
\define@key[psset]{pst-numeric}{regPlot}{\def\psk@regPlot{#1}}
\psset[pst-numeric]{numFormat=raw,numDigits=6,numTrim=false,numDecimal=.,
  sciStyle=times,matrixDelim=bracket,matrixAlign=decimal,colSep=1,rowSep=1.4,
  vector=column,label={},tensorLabels=true,regShow={eq r2 s},regPlot=fit}
% the settings as PostScript, for a dictionary on top of tx@NumericDict
\def\pstnum@setup{%
  /pstnum_format /\psk@numFormat\space def
  /pstnum_digits \psk@numDigits\space def
  /pstnum_trim \ifPst@numTrim true\else false\fi\space def
  /pstnum_dsep (\psk@numDecimal) def
  /pstnum_sci /\psk@sciStyle\space def
  /pstnum_font /\psNumericFont\space def
  /pstnum_fsize \psNumericFontSize\space def
  /pstnum_delim /\psk@matrixDelim\space def
  /pstnum_align /\psk@matrixAlign\space def
  /pstnum_colsep \psk@colSep\space def
  /pstnum_rowsep \psk@rowSep\space def
  /pstnum_vector /\psk@vector\space def
  /pstnum_label (\psk@numLabel) def
  /pstnum_tlabels \ifPst@tensorLabels true\else false\fi\space def
  /pstnum_regshow (\psk@regShow) def }
% #1 options, #2 PostScript drawing code run with the settings
\def\pstnum@draw#1#2{%
  \pscustom[#1]{\moveto(0,0)%
    \code{tx@NumericDict begin 40 dict begin \pstnum@setup #2 end end}}}
%
% \psNumericShow[options]{code} prints the result that the code leaves on
% the stack, at the current point, as a zero-width object: put it in a
% pspicture that leaves room for the text. With numFormat=raw (the default)
% a bignum is printed with all its digits, a complex number as a + bi, a
% matrix as [a b; c d]; the other formats round as asked.
\def\psNumericFont{Courier}
\def\psNumericFontSize{10}
\def\psNumericShow{\pst@object{psNumericShow}}
\def\psNumericShow@i#1{%
  \begingroup\use@par\pstnum@draw{}{#1 pstnum_textdraw}\endgroup\ignorespaces}
%
% \psNumericMatrix[options]{code} draws the result laid out: a matrix with
% delimiters and aligned columns, a list of numbers as a vector, a tensor of
% rank 3 or 4 as its slices; the left end of the math axis is at (0,0).
% The default format is auto.
\def\psNumericMatrix{\pst@object{psNumericMatrix}}
\def\psNumericMatrix@i#1{%
  \begingroup\use@par\pstnum@draw{}{/pstnum_font /\psNumericMatrixFont\space def #1 pstnum_mdraw}\endgroup\ignorespaces}
\def\psNumericMatrixFont{Times-Roman}
%
% \psNumericRegression[options]{code}: a regression result (from regress)
% as lines of text, the first baseline at (0,0); regShow picks the lines.
\def\psNumericRegression{\pst@object{psNumericRegression}}
\def\psNumericRegression@i#1{%
  \begingroup\use@par\pstnum@draw{}{/pstnum_font /\psNumericMatrixFont\space def #1 pstnum_regdraw}\endgroup\ignorespaces}
%
% \psNumericDef{name}{code}: runs the code once, in the dvips header, and
% stores its result as /name in tx@NumericDict, for every page. Use it for
% a result needed many times, such as a regression for \psNumericRegPlot or
% \psplot. Needs the latex + dvips route.
\def\psNumericDef#1#2{\special{! tx@NumericDict begin /#1 #2 def end}}
%
% \psNumericRegPlot[options]{name}{x0}{x1}: plots the regression stored as
% name (\psNumericDef) from x0 to x1: regPlot=fit (the curve), ci or pi
% (both limits of the confidence or prediction band), data, residuals or
% stdresiduals (as dots, at the data's x values).
\def\psNumericRegPlot{\pst@object{psNumericRegPlot}}
\def\psNumericRegPlot@i#1#2#3{%
  \begingroup\use@par
  \def\pstnum@tmp{fit}\ifx\psk@regPlot\pstnum@tmp
    \psplot{#2}{#3}{x tx@NumericDict begin #1 exch regeval pstnum_real end}\fi
  \def\pstnum@tmp{ci}\ifx\psk@regPlot\pstnum@tmp
    \psplot{#2}{#3}{x tx@NumericDict begin #1 exch regci 0 get pstnum_real end}%
    \psplot{#2}{#3}{x tx@NumericDict begin #1 exch regci 1 get pstnum_real end}\fi
  \def\pstnum@tmp{pi}\ifx\psk@regPlot\pstnum@tmp
    \psplot{#2}{#3}{x tx@NumericDict begin #1 exch regpi 0 get pstnum_real end}%
    \psplot{#2}{#3}{x tx@NumericDict begin #1 exch regpi 1 get pstnum_real end}\fi
  \def\pstnum@tmp{data}\ifx\psk@regPlot\pstnum@tmp
    \listplot[plotstyle=dots]{\pstNumeric{#1 /y get #1 pstnum_xy}}\fi
  \def\pstnum@tmp{residuals}\ifx\psk@regPlot\pstnum@tmp
    \listplot[plotstyle=dots]{\pstNumeric{#1 /resid get #1 pstnum_xy}}\fi
  \def\pstnum@tmp{stdresiduals}\ifx\psk@regPlot\pstnum@tmp
    \listplot[plotstyle=dots]{\pstNumeric{#1 /stdresid get #1 pstnum_xy}}\fi
  \endgroup\ignorespaces}
%
\catcode`\@=\PstAtCode\relax
\endinput
